\chapter{Introduction}\label{ch:intro}

\section{Motivation and problem}\label{sec:motivation}

An old photograph is rarely damaged in just one way. A print shot through a
soft lens on grainy film, scanned at low resolution and saved as a
compressed JPEG carries all of those defects, and each one hits a different
part of the spectrum. Blur removes high frequencies. Resampling throws away
detail above a resolution limit. Noise dominates where the signal is weak,
which in natural images means the high frequencies. JPEG quantises the mid
and high frequencies of each $8\times8$ block. Most restoration networks
treat the mixture as a single problem, and the strongest of them, such as
Restormer \cite{zamir2022restormer}, need tens of millions of parameters.

We ask whether a very small network can use this spectral structure
directly. It splits the image into frequency bands, and a learned gate
chooses between convolutional experts of different kernel sizes for every
band and every pixel. We call the network \fcsg{}.

Formally, we degrade a clean RGB image $f\in[0,1]^{H\times W\times 3}$ as
$x=\mathcal{D}(f;\theta)$ with $\theta=(\sigma_\text{blur}, s, \sigma_n, q)$.
$\mathcal{D}$ applies Gaussian blur, bicubic down- and up-sampling by $s$,
Gaussian noise with standard deviation $\sigma_n$, and JPEG at quality $q$.
We draw the parameters at random for each image and hide them from the
model. We want $\hat f=g_\phi(x)$ close to $f$ in PSNR, SSIM and LPIPS, with
fewer than five million parameters and fewer than ten GFLOPs for a
$256\times256$ input.

\section{Objectives and contributions}\label{sec:objectives}

We test one hypothesis. \emph{Apply a fixed frequency decomposition, then
route each band per pixel to the top two of several shared kernel-size
experts. The routing will change with content and degradation, and the
trained network will depend on it, all within the budget.} The work has four
objectives.
\begin{enumerate}
	\item[\textbf{O1}] Build \fcsg{} and verify it before training.
	\item[\textbf{O2}] Get trustworthy baselines by reproducing a published
	FFDNet result.
	\item[\textbf{O3}] Compare \fcsg{}, DnCNN and two FFDNet variants trained
	under one matched protocol, in quality, size and measured speed.
	\item[\textbf{O4}] Describe what the router learned.
\end{enumerate}

We contribute a 97,088-parameter network with its verification suite, a
matched composite-degradation benchmark, and a routing analysis that links
expert weights to the recorded degradation. We also report the outcome as it
is. \fcsg{} comes within 0.6\dB{} of DnCNN with 5.7 times fewer parameters,
but it is not faster on a GPU. We claim neither state-of-the-art quality nor
novelty, since frequency decomposition and expert routing both have a long
history (Chapter~\ref{ch:literature}).

Chapter~\ref{ch:literature} reviews related work and
Chapter~\ref{ch:method} describes the method. Chapter~\ref{ch:setup} covers
the experimental setup and verification. Chapter~\ref{ch:results} gives the
results and the routing analysis, and Chapter~\ref{ch:conclusion} concludes.
